% 本文件由 LaTeX 排版 Agent 根据有效 translation.json 自动生成。
% author=Hilary of Poitiers; work_id=3301; title=论大公会议

\chapter{论大公会议}

% structure: p0001 source=h1 target=omitted reason=page-title-already-rendered-by-parent

% structure: p0002 source=h2 target=section reason=relative-heading-rank; base=h2

\section{引言}

\Person{希尔拉里}并未参加公元358年春在\Place{安西拉}举行的主教会议，但他已了解该会议的决定，甚至了解了该会议的使节在\Place{希尔米乌姆}所隐瞒的绝罚。他看出这些决定标志着一个接近。那些\Person{尤西比安}派人士（他们对\Place{尼西亚}信经的唯一反对是并不理解它）对\Place{希尔米乌姆}的\OriginalTerm{亵渎文}{Blasphemia}所感到的恐惧，预示着未来会好转。同时，大多数东方主教是有意异端的。\Person{希尔拉里}自然对西方的主教职感到忧虑。\par

他已流亡约三年，并与西方主教们通信。从几个方面寄来的信件现已停止，他害怕主教们不愿写信给一个信念与他们不同的人。当他于358年底收到一封信时，他不胜欣喜，信不仅解释说他们沉默的无辜原因是不知道他的地址，而且还说他们坚持拒绝与\Person{撒图尔尼努斯}共融，并谴责了\OriginalTerm{亵渎文}{Blasphemia}。\par

359年初，他向他们派发了\WorkTitle{论主教会议}。这是一封双重的书信，写给西方主教们，但也包含针对东方人（他们无疑会及时收到这封信）的段落。\Person{希尔拉里}认识到，西方的正统信徒一直与东方的正统信徒保持距离，首先是因为对事件的无知，其次是因为对\OriginalTerm{同体}{\GreekText{ὁμοούσιος}}一词的误解，第三是因为当时普遍存在的猜疑情绪。这些事实决定了他书信的内容。\par

他以表达他在得知\Place{高卢}主教们谴责了臭名昭著的\Place{希尔米乌姆}信条时所感到的喜悦开始。他赞扬他们信仰的坚定。\par

他接着提到，他收到他们中某些人的请求，要他提供一份关于在东方所制定的诸信经的报告。他谦逊地答应了这一请求，恳求他的读者在读完整封信并掌握完整的论证之前不要批评他的信。他的目标始终是挫败异端者并帮助天主教徒。\par

在书信的第一部分或历史部分，他承诺提供一份自\Place{尼西亚}大公会议以来所制定的所有信经的抄本及解释，\par

\SourceRecord{Translated by E.W. Watson and L. Pullan.；Nicene and Post-Nicene Fathers, Second Series；Vol. 9.；Edited by Philip Schaff and Henry Wace.；Buffalo, NY: Christian Literature Publishing Co.,；1899.；<http://www.newadvent.org/fathers/3301.htm>.}
